\section{Frames and their maps}
\label{sec:A-frames}

The word ``frame'' carries the conjecture of Section~\ref{sec:A-intro}. This section fixes what it
means in the construction, which maps connect the frames, and what such maps can and cannot deliver.
The maps and their properties are \Established; the identification of frames with physical
observers of matter, light and tachyons is a \Reading.

\subsection{The field doublet and its state disk}
\label{sec:A-doublet}

Write the field as the real $3\times2$ matrix $F=[\,\mathbf E\mid\mathbf B\,]$ and form its Gram matrix
$G=F^{\mathsf T}F$ (\rndold{r04}). With $u=\tfrac12\operatorname{tr}G=\tfrac12(|\mathbf E|^2+|\mathbf B|^2)$ and
$s$ from \eqref{eq:A-s},
\begin{equation}
G=u\,\mathbf 1+\tfrac12\bigl(\operatorname{Re}s\;\sigma_z+\operatorname{Im}s\;\sigma_x\bigr),
\qquad \det G=|\mathbf E\times\mathbf B|^2,\qquad u^2-\det G=\tfrac14|s|^2 .
\label{eq:A-gram}
\end{equation}
For $u>0$ the normalised matrix $\rho=G/\operatorname{tr}G=\tfrac12(\mathbf 1+b_1\sigma_z+b_2\sigma_x)$
defines the \emph{state vector}
\begin{equation}
b=\frac{(\operatorname{Re}s,\operatorname{Im}s)}{2u},\qquad
|b|^2=1-\frac{\det G}{u^2}\le1 .
\label{eq:A-disk}
\end{equation}
The states fill a disk, not a sphere, because $G$ is real symmetric. Light-like fields ($s=0$) are
the centre $b=0$. The rim $|b|=1$ is reached exactly when $\mathbf E\parallel\mathbf B$. A boost leaves
$s$ fixed and changes $u$, so it moves $b$ along its own direction; for a non-null field it reaches
the rim in the frame where $\mathbf E\parallel\mathbf B$ and tends to the centre as the rapidity grows
(\rndold{r04}, r04.2). This is the Stokes formalism of polarisation optics, applied to the pair
$(\mathbf E,\mathbf B)$ instead of two polarisation components; the same three bases of the field
doublet appear in the literature \cite{A-Golat2024}.

\subsection{Convention F}
\label{sec:A-convention}

Since round \rndold{r22} every statement of the Companion carries the frame in which it holds. The
convention, completed in \rndold{r24} and corrected in \rndold{r28}, names four frames.
\begin{itemize}
\item \textbf{M}, the matter frame: the doublet state with eigenvalues $p_+=1/(1+x)$ in the compensator
  port and $p_-=x/(1+x)$ in the partner port, where $x=\bsym^2$.
\item \textbf{T}, the tachyon frame: the image of M under the tachyon map $z\mapsto iz$. It exchanges
  $p_+$ and $p_-$, sends $x\mapsto1/x$ and the frame factor $\gamma=(1+x)/(1-x)$ to $-\gamma$.
\item \textbf{P}, the photon eigenframe: the fixed point $b=0$, $p_+=p_-=\tfrac12$.
\item \textbf{R}, the reference frame: the pure state $|b|=1$, offset $\bsym=0$. It is the frame in
  which the Companion main text works; where older chapters say ``photon frame'' they mostly mean R.
\end{itemize}
A statement is \emph{frame-free} if M and T give the same value; otherwise it carries its frame. At P
a readout is evaluated separately and may be degenerate, and a degenerate value is not a global
statement. Quantities that depend on the state only through $|b|$ or the unordered pair of eigenvalues
(visibility, entropy) are frame-free; quantities that distinguish the two ports are matter-only
(\rndold{r22}, r22.1). The positions of M and T depend on the value of $\bsym$, which is inserted
(Section~\ref{sec:beta}); the convention itself does not.

\subsection{Three maps on the frame plane}
\label{sec:A-maps}

On the complex amplitude $z$, and in the Dixon algebra on its $\mathbb C$ factor (\rndold{r51}),
three involutive operations act. We write them as
\begin{equation}
\begin{aligned}
&D\colon z\mapsto iz &&\text{(tachyon map)},\\
&\vartheta\colon z\mapsto\bar z &&\text{(time reversal of the amplitudes)},\\
&J_B=D\circ\vartheta\colon z\mapsto i\bar z &&\text{(exchange of $\mathbf E$ and $\mathbf B$)}.
\end{aligned}
\label{eq:A-maps}
\end{equation}
On the frame plane $D$ is a rotation by $90^\circ$ and $\vartheta$ a reflection
(Figure~\ref{fig:A-frame-plane}). Strictly, $D$ squares to $-1$ on $z$; it is an involution on states,
where it acts as $\rho\mapsto\sigma_y\rho\sigma_y$. In the matrix picture $\vartheta$ is $F\mapsto F\sigma_z$,
$J_B$ is $F\mapsto F\sigma_x$, and $D$ is the rotation by $\pi/2$ in the column plane.

\begin{proposition}[Frame maps; v4.8.0-r06, \rndold{r04}, \rndold{r28}, \rnd{r07}]
\label{prop:A-maps}
\begin{enumerate}
\item[(a)] $s(Dz)=-s(z)$, $s(\vartheta z)=\overline{s(z)}$, and $s(J_Bz)=-\overline{s(z)}$. On the disk, $D$ sends
  $b\mapsto-b$, $\vartheta$ reverses $b_2$, and $J_B$ reverses $b_1$. All three keep $|b|$ and the null class.
\item[(b)] $D$ commutes with every $L\in SO(3,\mathbb C)$, that is, with all proper Lorentz transformations
  of the field.
\item[(c)] $\vartheta$ and $J_B$ commute with $L\in SO(3,\mathbb C)$ if and only if $L=\bar L$, that is, for
  spatial rotations only. A boost changes the form of the exchange rule.
\end{enumerate}
\end{proposition}

\begin{proof}
(a) $s$ is complex bilinear: $s(iz)=i^2s(z)$, $s(\bar z)=\overline{s(z)}$, and $s(i\bar z)=-\overline{s(z)}$. The
statements on $b$ follow from \eqref{eq:A-disk}, since $u$ is unchanged by all three maps. (b) $L$ is
complex linear, so $L(iz)=iLz$. (c) $L\vartheta z=L\bar z$ and $\vartheta Lz=\bar L\bar z$; these agree for all
$z$ iff $L=\bar L$. Real elements of $SO(3,\mathbb C)$ are the rotations. The same holds for $J_B=D\vartheta$
by (b).
\end{proof}

Three consequences need to be stated plainly, because earlier text of the Companion confused them.

\paragraph{The tachyon map is not time reversal.}
Both $D$ and $J_B$ reverse $\operatorname{Re}s$ and so send matter-type to tachyon-type configurations;
$\vartheta$ keeps $\operatorname{Re}s$. Convention~F as first written used $J_B$; round \rndold{r28}
corrected the tachyon map to $z\mapsto iz$, which in addition exchanges the compensator ports, and round \rndold{r51} showed that the charges $Q$ and $B-L$ and the chirality projectors of
the Dixon construction are fixed by $D$, so they are the same in M and T, whereas charge conjugation is
the time-reversal part $\vartheta$ \Established. Where earlier layers wrote ``frame T'' for charge conjugation,
time reversal was meant. The two maps must not be interchanged.

\paragraph{The tachyon map is electromagnetic duality, not a Wick rotation.}
With $z=\mathbf E+i\mathbf B$, $Dz$ is the duality rotation $\mathbf E\mapsto-\mathbf B$,
$\mathbf B\mapsto\mathbf E$, a map of field space and not of spacetime. By (b) it commutes with every
boost; a boost stays a boost with real rapidity. A remark in \rnd{r06} had asked whether $z\mapsto iz$
turns boosts into rotations and coincides with the Wick rotation $t\mapsto i\tau_E$; \rnd{r07} (Correction
r07.1) answered this in the negative \Established.

\paragraph{The names are a reading.}
That the orbit related by $D$ to a matter-type configuration is a physical tachyon, and that the null
class is light in the particle sense, is the working assignment \eqref{eq:A-assignment}. The maps
themselves are symmetries of the free Maxwell field (duality, Lorentz transformations) or a discrete
transformation of it (time reversal).

% F3: the frame plane and the state disk.  Tachyon map z -> iz (rotation by 90 deg) against time reversal z -> conj z
% (reflection) and J_B = i conj z.  2.5D: two stacked panels with shadows; the lower panel is the readout disk b.
% All positions are exact consequences of the maps: b(iz) = -b(z), b(conj z) = (b1,-b2), b(J_B z) = (-b1,b2)
% (Companion v4.9.0-r04.2(b),(c), v4.9.0-r28).
\begin{figure}[t]
\centering
\begin{tikzpicture}[
  font=\footnotesize,
  vec/.style={-{Stealth[length=2.2mm]}, line width=0.9pt},
  panel/.style={fill=white, draw=black!55},
  shadow/.style={fill=black!22},
]
% ---------------- upper panel: the frame plane ----------------
\begin{scope}
  \fill[shadow] (-2.95,-2.45) rectangle (3.05,2.65);
  \filldraw[panel] (-3.1,-2.3) rectangle (2.9,2.8);
  \node[anchor=north west, font=\footnotesize\bfseries, text=black!65] at (-3.05,2.75) {frame plane (the $\mathbb C$ factor)};
  \draw[-{Stealth[length=1.6mm]}, black!60] (-2.4,0) -- (2.5,0) node[below left] {$\operatorname{Re}$};
  \draw[-{Stealth[length=1.6mm]}, black!60] (0,-2.0) -- (0,2.3) node[below right] {$\operatorname{Im}$};
  \draw[black!25] (0,0) circle (1.8);
  % z at 30 deg, iz at 120 deg, conj z at -30 deg, J_B z = i conj z at 60 deg
  \draw[vec, black] (0,0) -- (30:1.8) node[right] {$z$};
  \draw[vec, black, dashed] (0,0) -- (120:1.8) node[above left] {$iz$};
  \draw[vec, black, densely dotted] (0,0) -- (-30:1.8) node[right] {$\bar z$};
  \draw[vec, black!60, dash dot] (0,0) -- (60:1.8) node[above right] {$J_Bz=i\bar z$};
  \draw[-{Stealth[length=1.5mm]}, black] (30:0.75) arc (30:120:0.75);
  \node at (75:1.02) {$90^\circ$};
  \draw[{Stealth[length=1.5mm]}-{Stealth[length=1.5mm]}, black!70] (30:1.35) arc (30:-30:1.35);
  \node[right] at (8:1.38) {\scriptsize mirror};
\end{scope}
% ---------------- lower panel: the state disk ----------------
\begin{scope}[xshift=7.2cm, yshift=-0.9cm]
  \fill[shadow] (-2.85,-2.45) rectangle (3.15,2.65);
  \filldraw[panel] (-3.0,-2.3) rectangle (3.0,2.8);
  \node[anchor=north west, font=\footnotesize\bfseries, text=black!65] at (-2.95,2.75) {state disk $b=(\operatorname{Re}s,\operatorname{Im}s)/2u$};
  \fill[black!6] (0,0) circle (1.8);
  \draw[black!60] (0,0) circle (1.8);
  \draw[black!45] (-2.1,0) -- (2.1,0) node[below left] {$b_1$};
  \draw[black!45] (0,-2.0) -- (0,2.15) node[below right] {$b_2$};
  \node[font=\scriptsize, text=black!70, align=center] at (1.45,-2.08) {$b_1>0$: matter};
  \node[font=\scriptsize, text=black!70, align=center] at (-1.45,-2.08) {$b_1<0$: tachyon};
  % P at centre; a boost moves b radially to the rim (frame R)
  \fill (0,0) circle (1.6pt) node[above left=-1pt] {$P$};
  \draw[-{Stealth[length=1.5mm]}, black!70] (1.2,0.709) -- (1.503,0.888);
  \draw[fill=white] (1.550,0.916) circle (1.6pt) node[above=1pt] {$R$};
  % sample state b and its images
  \fill (1.1,0.65) circle (1.5pt) node[above left=0pt and -2pt] {$b$};
  \fill (-1.1,-0.65) circle (1.5pt) node[left=1pt] {$-b$};
  \fill (1.1,-0.65) circle (1.5pt) node[below=1pt] {$(b_1,-b_2)$};
  \fill (-1.1,0.65) circle (1.5pt) node[above=1pt] {$(-b_1,b_2)$};
  \draw[-{Stealth[length=1.6mm]}, dashed] (1.03,0.61) -- (-1.03,-0.61);
  \draw[-{Stealth[length=1.6mm]}, densely dotted] (1.1,0.57) -- (1.1,-0.57);
  \draw[-{Stealth[length=1.6mm]}, dash dot, black!60] (1.02,0.65) -- (-1.02,0.65);
\end{scope}
% connector between the panels
\draw[black!45, densely dashed, -{Stealth[length=1.6mm]}] (2.95,0.6) to[out=0,in=180] (4.15,0.1);
\end{tikzpicture}
\caption{The frame maps on the frame plane (left) and on the readout disk of the field doublet (right).
The tachyon map $z\mapsto iz$ (dashed) is a rotation by $90^\circ$; on the disk it is the point
reflection $b\mapsto-b$. Time reversal of the amplitudes $z\mapsto\bar z$ (dotted) is a mirror image;
on the disk it reverses $b_2$. Their composite $J_B z=i\bar z$ (dash-dotted) reverses $b_1$. Light-like
fields sit at the centre $P$ in every frame; a boost moves a non-null state along its own direction
and reaches the rim $|b|=1$, the frame $R$ with $\mathbf E\parallel\mathbf B$, at a finite rapidity.
\Established{} for the maps and their action \cite{QSC} (\rndold{r04}, \rndold{r28}, \rndold{r51}).}
\label{fig:A-frame-plane}
\end{figure}


\subsection{Exactly two classes of fields}
\label{sec:A-classes}

If one divides out not only the Lorentz group and duality but also the overall scale, the sector
structure collapses to a single predicate.

\begin{proposition}[Two classes; \rnd{r21}, r21.8]
\label{prop:A-classes}
Let $SO(3,\mathbb C)$ act on $z$ and let $\lambda\in\mathbb C\setminus\{0\}$ act by $z\mapsto\lambda z$ (duality
phase and dilation).
\begin{enumerate}
\item[(a)] All non-null fields form one orbit: for $s(z_1),s(z_2)\neq0$ there are $\lambda$ and $L$ with
  $z_2=\lambda Lz_1$.
\item[(b)] The non-zero null fields form one orbit, already of $SO(3,\mathbb C)$ alone.
\item[(c)] Each class has a distinguished representative. A non-null field is taken by a finite boost to
  the frame with $\mathbf E\parallel\mathbf B$, where $\phi:=|b|=1$; this is R. A null field has $\phi=0$ in
  every frame; this is P. On the first class boosts reach every $\phi\in(0,1]$; on the second only $0$.
\item[(d)] The rapidity $\eta_R$ of the boost from a field to R satisfies $\phi=\operatorname{sech}2\eta_R$.
\end{enumerate}
\end{proposition}

\begin{proof}
(a) Choose $\lambda$ with $\lambda^2s(z_1)=s(z_2)$; then $\lambda z_1$ and $z_2$ have the same $s$ and lie on
one $SO(3,\mathbb C)$-orbit by Proposition~\ref{prop:A-orbits}(a). (b) is Proposition~\ref{prop:A-orbits}(a)
for the null cone. (c) and (d): $s$ is fixed and $\phi=|s|/2u$; along the boost line $u$ takes every value
in $[|s|/2,\infty)$, with the minimum where $\det G=0$, i.e.\ $\mathbf E\parallel\mathbf B$; for $s=0$,
$\phi=0$ identically. The relation (d) was checked numerically on 1000 fields in \rnd{r21}.
\end{proof}

Below every frame-dependent readout there is therefore one frame-free predicate: whether a field is
null or not \Established. The port shares, the offset and the value of $\phi$ are readouts of a
frame above it. This is the precise sense in which matter and tachyon are ``one'' and light is
``another'': the first two lie in one class, light in the other, and no finite boost connects them.

\subsection{The collapse state, typed}
\label{sec:A-collapse}

The oldest formulation of the conjecture calls the common object the \emph{collapse state}
$z^-=0$ (\texttt{thm:unity}). Round \rnd{r22} typed all 93 occurrences of $z^-$ in the main text against
Convention~F \Reading. As a point, $z^-=0$ is in about two thirds of the places the zero of the offset,
that is R. The role of the frame-invariant object, which the same theorem ascribes to $z^-=0$, belongs to
P, and the main text places P, without naming it so, at the other end of its coordinate, $z^-\to\infty$
at $\beta\to1$, the horizon. The coordinate thus runs from R to P along the line
$\phi=\operatorname{sech}2\eta_R$ of Proposition~\ref{prop:A-classes}(d). The null class, the answer to
the question whether a field is null at the lowest level, is the horizon end of the main-text coordinate,
not its origin.

The same split resolves where the main-text ``cascade'' leads (\rnd{r23}). Two cascades carry one name.
The compensator cascade ends at R: every image $Cz$ with $\mathbf E+\mathbf B\neq0$ has $\phi=1$
(Section~\ref{sec:A-field-example}). The transport of the null axis keeps null fields null and so carries
P without leading there. No finite cascade in the source leads to P; only limits (boosts, dephasing over
many frequencies, weak coupling to a reservoir) or, for an instant, free time evolution reach it.

\subsection{Frame maps give equalities, not values}
\label{sec:A-equalities}

A symmetry between frames can show that two frames agree. It cannot supply the number on which they
agree. The clearest example in the Companion concerns gravity.

Dorau and Much derived the semiclassical Einstein equations from quantum relative entropy
\cite{A-DorauMuch2026}, a quantum field theoretic form of Jacobson's argument
\cite{A-Jacobson1995}. Near a point, with a local Rindler horizon, the relative entropy of a coherent
excitation equals an energy flux; with the Raychaudhuri equation and the area law $S_{\rm rel}=\delta A/4$,
which is assumed, one obtains $G_{ab}+\Lambda g_{ab}=8\pi\langle{:}T_{ab}{:}\rangle$, with $\Lambda$ an
integration constant. Round \rnd{r06} recorded this; round \rnd{r07} mapped the chain onto the frames
\Exploration{} (not blind):

\begin{proposition}[Frame-blind gravity for the Maxwell field; \rnd{r07}, r07.2]
\label{prop:A-frameblind}
For the free Maxwell field the relative-entropy chain gives in frame T the same Einstein equation as in
frame M, with the same coupling, sign and inputs.
\end{proposition}

\begin{proof}[Sketch]
The duality rotation is a symmetry of the free Maxwell field: it maps the algebra of the Rindler wedge
onto itself and fixes the vacuum. It commutes with boosts (Proposition~\ref{prop:A-maps}(b)), hence with
the modular group of the vacuum. Relative entropy is invariant under such an automorphism. The
energy--momentum tensor is Hermitian in $z$ (energy density $\tfrac12|z|^2$, Poynting vector
$(\bar z\times z)/2i$), so the null flux $T_{ab}\xi^a\xi^b$ is unchanged, while the bilinear invariants
$|\mathbf E|^2-|\mathbf B|^2$ and $\mathbf E\cdot\mathbf B$ change sign. The commutation and the invariance of
the tensor were checked numerically on 2000 random fields and boosts (script
\texttt{k217\_dualitaet\_gravitation.py}).
\end{proof}

Gravity couples to the Hermitian tensor; the frames differ in the bilinear invariants. For the Maxwell
field gravity is therefore frame-blind. This is textbook duality invariance in a new bookkeeping, not
new physics. Its limits are explicit: for massive fields the tachyon image of the action has no stable
Hadamard ground state, and the premises of the chain fail; for interacting fields duality is in general
not a symmetry.

The more important lesson is structural (\rnd{r07}, r07.4). The frame map is a symmetry that fixes the
reference state, and every link of the chain is invariant under it; such an argument yields equalities
between frames and nothing about the shared value. The prefactor of the area law is an absolute number.
The two horizon countings the Companion uses are each calibrated to the area law, and their ratio
$\gamma_L=\gamma_{\rm BI}/\gamma_{\rm string}=1.864700$ is independent of the normalisation, which cancels.
Every quantity the Companion builds from the two countings, including $\bsym$ (Section~\ref{sec:beta}), is exactly the
part in which the normalisation drops out. Neither Newton's constant nor $\Lambda$ follows
\Established.

\subsection{Orientation of the octant}
\label{sec:A-orientation}

One further property of the maps is used in Section~\ref{sec:B-koide}. On the Lorentz-invariant doublet
$\chi=(\sqrt{s_+},\sqrt{s_-})$ built from the helicity parts $F_\pm=\mathbf E\pm i\mathbf B$ of a field with
complex amplitudes (\rnd{r16}), the eigenstates of the frame maps form octant triples of three states with
pairwise overlap $\tfrac12$; the eight triples split into two orbits, one for each sign of the Bargmann phase
$\pm\pi/4$. The pre-registered test H118 (\rnd{r31}, partly not blind) found that the three frame maps, acting
linearly on the complex amplitudes, preserve this orientation, and so does parity, which exchanges $s_+$
and $s_-$. Time reversal of the complex field amplitudes, $\mathbf E\mapsto\bar{\mathbf E}$,
$\mathbf B\mapsto-\bar{\mathbf B}$, acts on $\chi$ as complex conjugation and reverses it. None of nine candidates listed beforehand
selects one orientation. The sign of the octant step is therefore a time orientation of the readout and is
not fixed by the construction \Open. Section~\ref{sec:B-koide} returns to this.

\subsection{Summary}
\label{sec:A-frames-summary}

\begin{table}[h]
\centering\small
\begin{tabularx}{\textwidth}{@{}>{\raggedright\arraybackslash}X l l@{}}
\toprule
statement & status & source\\
\midrule
natural norm $N(q)=-t^2+s(z)$; Lorentz action on $z$ (Reading~B) & \Established & v4.8.0-r06\\
compensator $C=\tfrac12(1+J)$ is an orthogonal projection & \Established & v4.8.0-r04\\
orbits and sector rule by the sign of $s$ & \Established & v4.8.0-r06, r07\\
names matter / light / tachyon for $s>0$, $=0$, $<0$ & \Reading & v4.8.0-r07\\
image of $C$ never light-like; lies at R & \Established & \rnd{r23}\\
$z\mapsto iz$ is duality, commutes with boosts; not a Wick rotation & \Established & \rnd{r07}\\
tachyon map $\neq$ time reversal; charges frame-free between M and T & \Established & \rndold{r28}, \rndold{r51}\\
two classes of non-zero fields, representatives P and R & \Established & \rnd{r21}\\
collapse state: point R, invariant-object role P & \Reading & \rnd{r22}\\
gravity frame-blind for the free Maxwell field & \Exploration{} (not blind) & \rnd{r07}\\
values (area-law prefactor, $G$, $\Lambda$) from frame maps & not obtainable & \rnd{r07}\\
orientation of the octant & \Open & \rnd{r31}\\
physical unity of matter, light and tachyon & \Open{} (conjecture) & \texttt{thm:unity}\\
\bottomrule
\end{tabularx}
\caption{Status of the results of Sections~\ref{sec:A-carrier} and~\ref{sec:A-frames}.}
\label{tab:A-summary}
\end{table}
